\chapter{Work, Energy, and Power}

\mathTopics{\pgRef{chap:vectors}}

\section{Energy}

In Newtonian mechanics, the kinetic energy \(E_\mathrm{K}\) is given by
\begin{equation} \label{eq:linear-kinetic-energy}
    E_\mathrm{K} = \frac{1}{2} m v^2,
\end{equation}
where \(m\) is the mass and \(v\) is the velocity.

The equation for \textbf{potential energy}~\(U\) is
\begin{equation} \label{eq:potential-energy}
    U = - \int \vecb{F} \cdot \odif{\vecb{r}}.
\end{equation}

\subsection{Gravitational Potential Energy}

Assuming a constant gravitational field,
the difference in gravitational potential energy \(U_\mathrm{g}\) is
\begin{equation}
    \begin{aligned}
        U_\mathrm{g} &= - \int - m g \odif{y} \\
        \Delta{U_\mathrm{g}} &= m g \Delta{y}.
    \end{aligned}
\end{equation}

If the gravitational field is from a sphere,
we use Newton's law of gravity [\pgRef{eq:newtons-law-of-gravity}]:
\begin{equation}
    \begin{aligned}
        U_\mathrm{g} &= - \int - G \frac{m_1 m_2}{r^2} \odif{r} \\
        \Delta{U_\mathrm{g}} &= - G \frac{m_1 m_2}{\Delta{r}}.
    \end{aligned}
\end{equation}

\subsection{Elastic Potential Energy}

The \textbf{elastic potential energy} for a general spring with force
function \(F(x)\) is
\begin{equation} \label{eq:general-elastic-potential-energy}
    U_\mathrm{e} = - \int F(x) \odif{x}
\end{equation}

For ideal springs, we use Hooke's law [\pgRef{eq:hookes-law}] to find
\begin{equation}
    \begin{aligned}
        \Delta{U_\mathrm{e}} &= - \int - k x \odif{x} \\
        &= \frac{1}{2} k (\Delta{x})^2
    \end{aligned}.
\end{equation}

The derivation from \pgRef{eq:general-elastic-potential-energy} can
also be used with non-linear springs.
For example, for a spring with force function \(F(x) = -k x^\frac{4}{3}\),
\begin{equation}
    \begin{aligned}
        U_\mathrm{e} &= - \int - k x^\frac{4}{3} \odif{x} \\
        \Delta{U_\mathrm{e}} &= \frac{3}{7} k x^\frac{7}{3}.
    \end{aligned}
\end{equation}

\section{Work}

\textbf{Conservative forces} are ones where the \textbf{work} done is
\textbf{path-independent},
i.e., work along a \textbf{closed loop} is 0.

The general formula for work \(W\) is
\begin{equation} \label{eq:work}
    \begin{aligned}
        W &= \int_C \vecb{F}(r) \cdot \odif{\vecb{r}} \\
        &= \int_C \vecb{F}(t) \cdot \vecb{r}'(t) \odif{t} \\
    \end{aligned},
\end{equation}
where \(\vecb{r}\) is the displacement.

For a straight path and constant force,
the formula simplifies down to
\begin{equation*}
    W = \vecb{F} \cdot \vecb{r}.
\end{equation*}

Compare the equation for work [\pgRef{eq:work}],
with the equation for potential energy [\pgRef{eq:potential-energy}].
Therefore, for conservative forces, the work~\(W\) is given by
\begin{equation}
    W = - \Delta{U},
\end{equation}
where \(U\) is the potential energy.
Be careful about the signs for work.
If positive work is done \emph{by} the system, it \emph{loses} energy.
If positive external work is done \emph{upon} the system, it gains energy.

\section{Power}

Power is work done per unit of time.

The average power \(P_\mathrm{avg}\) is given by
\begin{equation}
    P_\mathrm{avg} = \frac{W}{\Delta{t}} = \frac{\Delta{E}}{\Delta{t}},
\end{equation}
where \(W\) is work done, \(\Delta{t}\) is the change in time,
and \(\Delta{E}\) is the change in energy.

The instantaneous power \(P\) is given by
\begin{equation}
    P = \odv{E}{t}.
\end{equation}
